\chapter{Introduction and Background}

\section{Introduction}

In recent years, the field of computer vision has witnessed significant advancements in
segmentation tasks, particularly with the advent of convolutional neural networks (CNNs) and
transformer-based architectures. These developments have propelled instance segmentation to new
heights, enabling models to delineate object boundaries with remarkable accuracy. However,
traditional segmentation methods often fall short when confronted with complex, implicit queries
that necessitate higher-order reasoning.

Addressing this gap, the task of reasoning segmentation has emerged, aiming to generate
segmentation masks in response to intricate textual instructions. While initial research
has primarily focused on image-based reasoning segmentation, the dynamic nature of video
introduces additional challenges and opportunities. Videos encapsulate temporal information,
object motion, and evolving contexts, making them inherently more complex than static images.
Consequently, extending reasoning segmentation to the video domain requires novel methodologies
that can effectively harness both spatial and temporal dimensions.

Recent endeavors have introduced multimodal large language models (MLLMs) to tackle reasoning segmentation
in videos. Models such as VISA \cite{yan_visa_2024}, VideoLISA \cite{bai_one_2024}, and ViLLa \cite{zheng_villa_2025}
have demonstrated the potential of integrating world knowledge and temporal understanding to produce temporally
consistent segmentation masks based on language instructions. These models leverage advanced techniques such as
temporal context aggregation, sparse dense sampling, and token-driven keyframe selection to address the unique
challenges posed by video data. Despite their promising results, several limitations persist, including issues
with model generalization, handling of occlusions, and the need for extensive computational resources.

\section{Related Work}

\subsection{Multi-Modal Large Language Models}

Multi-Modal Large Language Models (MLLMs) have become a major research focus due to their ability to
integrate and reason across multiple modalities, such as text, image, and video data. These models
combine the power of Large Language Models (LLMs) with sophisticated vision encoders to enable complex
multimodal tasks such as image captioning, video question answering, and video segmentation. The recent
advancements in this field are largely driven by innovative architectures that combine cross-attention mechanisms,
vision encoders, and large-scale datasets to improve model performance across various tasks.

In many MLLMs, the integration of vision and language is achieved by using cross-attention mechanisms to
dynamically fuse visual features with textual information. For instance, Flamingo \cite{alayrac_flamingo_2022}
uses cross-attention to enable the model to focus on relevant visual contexts when generating text. Other models,
such as BLIP-2 \cite{li_blip-2_2023} and mPLUG-OWL \cite{ye_mplug-owl_2024}, encode visual inputs through pre-trained
vision encoders (such as CLIP \cite{radford_learning_2021}) and integrate these with text embeddings from LLMs,
resulting in rich multimodal representations. OTTER \cite{huang_otter_2025}, LLaVA \cite{li_llava-onevision_2024},
and MiniGPT-4 \cite{zhu_minigpt-4_2023} further enhance these capabilities by leveraging instruction-following and
few-shot learning approaches, where they fine-tune the models on specialized datasets for improved performance.

The InternVideo2 \cite{wang_internvideo2_2024} model contributes significantly to video processing within the
MLLM landscape by utilizing VideoMAEv2 \cite{wang_videomae_2023} as a video encoder. This model incorporates
depth-wise 3D convolutions for fusing spatial and temporal information across frames, which is crucial for
tasks such as segmentation. Additionally, InternVideo2 introduces more efficient methods of handling video frames
by focusing on selective attention, reducing the computational burden while maintaining high-quality outputs.

LLaVA \cite{li_llava-onevision_2024}, a notable model in the multimodal landscape, leverages both a powerful
LLM (Qwen-2) \cite{yang_qwen2_2024} and a robust vision encoder (SigLIP) \cite{zhai_sigmoid_2023}. LLaVA's architecture
is designed to scale effectively, supporting large datasets and powerful model sizes while maintaining strong
language capabilities. The model uses a 2-layer MLP projector to transform visual features into a space compatible with
the LLM, enabling effective cross-modal interaction. LLaVA's Higher AnyRes strategy addresses resolution and token number
challenges by adjusting image resolution and token configurations depending on the task, such as single-image, multi-image,
or video understanding. This strategy allows the model to handle varying visual complexities while maintaining computational
efficiency. The approach ensures the model can generalize across different scenarios, making it versatile for both images and videos.

The development of these models illustrates the potential of MLLMs to tackle complex multimodal tasks, with a specific
focus on improving video understanding. Techniques such as advanced tokenization strategies, efficient visual representations,
and the integration of deep temporal understanding are critical steps forward in the evolution of multimodal reasoning
systems. As these models continue to scale, their ability to perform tasks such as video segmentation, especially in
long-duration and high-complexity scenarios, will become increasingly robust, pushing the boundaries of what is possible
in multimodal AI.

\subsection{Video Object Segmentation}

Video Object Segmentation (VOS) is a crucial task in computer vision that focuses on the identification and
tracking of objects across video frames. It involves both segmenting objects in individual frames and ensuring
the continuity of the segmentation over time, making it particularly challenging due to the need to handle
temporal dependencies, varying object appearances, and motion dynamics. Video segmentation tasks can be broadly
categorized into Video Instance Segmentation (category-based VOS) and Referring Video Object Segmentation (text-based VOS),
each presenting unique challenges and requiring different strategies for effective solutions.

Video Instance Segmentation (VIS) aims to detect, segment, and track object instances across video frames,
typically based on predefined object categories. It has gained significant attention due to its complexity,
which involves modeling both spatial and temporal relationships of objects. Various approaches, such as tracking
by detection, masked attention, object token association, and contrastive memory, have been explored to enhance
VIS models. For instance, IDOL \cite{wu_defense_2022}, built on Deformable-DETR \cite{zhu_deformable_2021},
incorporates contrastive learning to generate unique embeddings for object associations, while DVIS \cite{zhang_dvis_2023}
decouples VIS into stages including image segmentation, tracking, and refinement. These advancements make VIS models
more robust and applicable across different scenarios.

Referring Video Object Segmentation (RVOS) introduces another level of complexity by segmenting target
objects mentioned in natural language expressions within a video. Unlike VOS, which focuses on tracking objects
based on predefined categories or initial frames, RVOS must handle both the action and appearance of the
referred object across multiple frames, making it more challenging. Methods such as URVOS \cite{vedaldi_urvos_2020}
and RefVOS \cite{bellver_refvos_2020} use cross-modal attention mechanisms to perform frame-by-frame segmentation,
while recent models such as ReferFormer \cite{zhao_learning_2023} and MTTR \cite{botach_end--end_2022} adopt
DETR-like structures to simplify the referring pipeline. Despite these advances, RVOS models still struggle
with understanding temporal dependencies and reasoning across long video sequences, especially in cases where
motion and contextual knowledge are crucial.

The Segment Anything Model 2 (SAM 2) \cite{ravi_sam_2024} offers a promising advancement in RVOS by
introducing a more robust memory attention mechanism that integrates both the visual information from past
frames and textual descriptions. SAM 2 conditions the current frame's features on previous frame features,
past predictions, and any new prompts, making it highly suited for RVOS tasks that involve complex referencing.

Through its memory encoder and memory bank, SAM 2 stores and retrieves information from past frames,
helping it maintain consistent object segmentation even across long sequences. This allows SAM 2 to better
handle scenarios where the referred object undergoes significant motion or changes in appearance over time.
Additionally, by leveraging both spatial and temporal memories, SAM 2 enhances its ability to segment and
track objects over long video sequences, making it highly effective for both simple and complex referring expressions in RVOS.

SAM 2's efficient integration of temporal reasoning alongside its advanced memory mechanisms sets it apart
from earlier RVOS methods, bringing it closer to handling more complex, real-world scenarios where reasoning
and inference are essential.

\subsection{Reasoning Segmentation}

Reasoning segmentation extends traditional visual recognition tasks by enabling models to interpret complex,
implicit textual queries and generate corresponding segmentation masks. This capability is crucial for
applications requiring nuanced understanding and interaction with visual content.

In the realm of image segmentation, models such as LISA \cite{lai_lisa_2024} have pioneered the integration
of large language models (LLMs) with segmentation frameworks. LISA employs a special $\langle$SEG$\rangle$ token to prompt
segmentation models, allowing them to produce masks in response to complex textual instructions. This approach
facilitates handling intricate reasoning tasks and demonstrates robust zero-shot capabilities.

Building upon this foundation, GSVA \cite{xia_gsva_2024} (Generalized Segmentation Vision Assistant)
introduces enhancements to address challenges in generalized referring expression segmentation (GRES).
GSVA extends the use of the $\langle$SEG$\rangle$ token to support multiple mask references simultaneously and introduces
a $\langle$REJ$\rangle$ token to explicitly reject null targets. These innovations enable GSVA to effectively manage
multiple-target and empty-target scenarios, outperforming previous methods on benchmarks such as gRefCOCO.

Extending reasoning segmentation to the video domain introduces additional complexities due to the
temporal dimension. Models must not only understand and segment objects but also track them across
frames in response to dynamic and evolving scenes.

VISA (Video-based Large Language Instructed Segmentation Assistant) \cite{yan_visa_2024} leverages
the capabilities of multimodal large language models to perform reasoning video object segmentation.
By integrating a mask decoder with world knowledge reasoning, VISA can generate segmentation masks for
objects based on complex textual prompts, facilitating tasks that require both visual understanding and temporal coherence.

VideoLISA \cite{bai_one_2024} builds upon the foundation of LISA by incorporating temporal dynamics
into the reasoning segmentation process. Utilizing a Sparse Dense Sampling strategy and a One-Token-Seg-All
approach with a specially designed $\langle$TRK$\rangle$ token, VideoLISA enables temporally consistent segmentation
across video frames. This design allows for effective tracking and segmentation of objects over time,
addressing the challenges posed by dynamic video content.

ViLLa (Video Reasoning Segmentation with Large Language Model) \cite{zheng_villa_2025} introduces
several innovations to enhance video reasoning segmentation. Its architecture includes a Key Segment
Extractor to identify relevant video segments, a Context Synthesizer to integrate visual and textual
information, and a Hierarchical Temporal Synchronizer to model multi-scale temporal dependencies. These
components collectively enable ViLLa to handle complex video scenarios involving multiple objects, occlusions,
and rapid motion, outperforming previous models on various benchmarks.

More recent advances have further pushed the boundaries of video reasoning segmentation (VRS). VRS-HQ \cite{gong_vrs_hq_2025}
introduces hierarchical temporal tokens, employing frame-level $\langle$SEG$\rangle$ and temporal-level $\langle$TAK$\rangle$ tokens
that leverage the MLLM's autoregressive learning to capture both local and global temporal information. Their
Token-driven Keyframe Selection (TKS) mechanism filters keyframes based on SAM 2's occlusion scores during inference,
achieving state-of-the-art performance on ReVOS with improvements of 5.9\% to 12.5\% over VISA across different subsets.

A paradigm shift has emerged with reinforcement learning-based approaches to VRS. Veason-R1 \cite{gong_veason_r1_2025}
is the first to apply Group Relative Policy Optimization (GRPO) augmented with Chain-of-Thought initialization,
enabling structured reasoning trajectories that bridge video-level semantics with frame-level spatial grounding.
This approach achieves significant improvements on ReasonVOS while exhibiting enhanced robustness to hallucinations.
Similarly, VideoSeg-R1 \cite{xu_videoseg_r1_2025} introduces a decoupled architecture that formulates VRS as joint
referring image segmentation and video mask propagation, incorporating a hierarchical text-guided frame sampler
and task difficulty-aware mechanisms for adaptive reasoning.

Notably, DecAF \cite{han_decaf_2025} demonstrates that competitive performance can be achieved without model retraining
through decomposed attention fusion. By extracting attention maps via rollout mechanisms and applying contrastive
object-background fusion with attention-guided SAM 2 prompting, DecAF achieves results comparable to training-based
methods in a training-free manner. This finding suggests that the rich visual-linguistic alignments learned during
pretraining may be sufficient for many VRS scenarios.

Additionally, DeSa2VA \cite{dang_desa2va_2025} addresses the feature entanglement problem observed in prior approaches
by proposing a decoupling-enhanced prompting scheme. Through linear projection that disentangles hidden states into
distinct textual and visual feature subspaces, combined with triple supervision from predicted text masks, visual masks,
and ground-truth annotations, DeSa2VA achieves state-of-the-art performance across diverse segmentation and grounding tasks.

These recent developments highlight several important trends in video reasoning segmentation: the transition from
supervised fine-tuning to reinforcement learning paradigms, the emergence of hierarchical token designs for better
temporal modeling, the viability of training-free approaches, and the importance of feature decoupling for improved
segmentation quality. These advances collectively address the core challenges of temporal consistency, computational
efficiency, and complex reasoning that remain central to the field.

